\documentclass[11pt,a4paper]{article}
\usepackage[utf8]{inputenc}
\usepackage[T1]{fontenc}
\usepackage[spanish,es-minimal,es-noshorthands]{babel}
\usepackage{lmodern}
\usepackage[sfdefault,lining]{FiraSans}
\usepackage[a4paper,top=2.3cm,bottom=2.2cm,left=1.8cm,right=1.8cm,headheight=1.3cm,headsep=0.4cm,footskip=1cm]{geometry}
\usepackage{amsmath,amssymb}
\usepackage{icomma}
\usepackage{graphicx}
\usepackage[table]{xcolor}
\usepackage{tikz}
\usepackage{fontawesome5}
\usepackage{ccicons}
\usepackage{fancyhdr}
\usepackage{lastpage}
\usepackage{tasks}
\usepackage{needspace}
\usepackage[most]{tcolorbox}
\usepackage{array}
\usepackage[hidelinks]{hyperref}
\AtBeginDocument{\hypersetup{pdftitle={Soluciones de la ficha de la unidad 12: Probabilidad. 3º ESO},pdfsubject={Matemáticas 3º ESO},pdfkeywords={probabilidad, experimento aleatorio, espacio muestral, sucesos, ley de Laplace, ley de los grandes números, frecuencia relativa, diagrama de árbol, reemplazamiento}}}

\definecolor{navy}{HTML}{1B2748}
\definecolor{accent}{HTML}{E8603A}
\definecolor{teal}{HTML}{1F9E90}
\definecolor{soft}{HTML}{EEF1F8}
\definecolor{muted}{HTML}{5D6685}

\setlength{\parindent}{0pt}
\setlength{\parskip}{3pt}
\settasks{label=\alph*), label-format=\bfseries\color{navy}, label-width=1.3em, item-indent=1.8em, column-sep=0.6em, after-item-skip=5pt, before-skip=3pt, after-skip=4pt}

% ---------- capçalera i peu ----------
\newcommand{\llicencia}{\href{https://creativecommons.org/licenses/by-nc-sa/4.0/deed.es}{\ccbyncsa}\ \ Licencia Creative Commons Atribución-NoComercial-CompartirIgual 4.0 Internacional (CC BY-NC-SA 4.0)}
\pagestyle{fancy}
\fancyhf{}
\fancyhead[L]{\raisebox{-0.35\height}{\includegraphics[height=0.95cm]{logo-gva.png}}}
\fancyhead[C]{\small\color{muted}Soluciones de la ficha · Unidad 12: Probabilidad · 3\textsuperscript{o} ESO}
\fancyhead[R]{\raisebox{-0.35\height}{\includegraphics[height=0.95cm]{logo-vermellar.png}}}
\fancyfoot[L]{\scriptsize\color{muted}\llicencia}
\fancyfoot[R]{\small\color{muted}\thepage}
\renewcommand{\headrulewidth}{0.4pt}
\renewcommand{\headrule}{\color{muted!40}\hrule height \headrulewidth}
\fancypagestyle{primera}{\fancyhf{}\renewcommand{\headrulewidth}{0pt}
  \fancyfoot[L]{\scriptsize\color{muted}\llicencia}
  \fancyfoot[R]{\small\color{muted}\thepage}}

% ---------- nivells i exercicis ----------
\newcommand{\estrelles}[1]{{\color{accent}\ifcase#1\or\faStar\faStar[regular]\faStar[regular]\or\faStar\faStar\faStar[regular]\or\faStar\faStar\faStar\fi}}
\newcounter{exer}
\newcommand{\ex}[1]{\par\needspace{4\baselineskip}\medskip\stepcounter{exer}\noindent{\color{navy}\bfseries\theexer.}\ \ {\small\estrelles{#1}}\ \ }

\newcommand{\seccio}[2]{\par\needspace{10\baselineskip}\bigskip
  \begin{tcolorbox}[colback=navy,colframe=navy,coltext=white,arc=2mm,boxrule=0pt,left=3mm,right=3mm,top=1.2mm,bottom=1.2mm]
  {\large\bfseries #1\quad #2}
  \end{tcolorbox}\vspace{-2pt}}

\newcolumntype{C}[1]{>{\centering\arraybackslash}p{#1}}
\usepackage{pifont}
\usetikzlibrary{trees}

% quadre de recordatori (només a la fitxa d'exercicis)
\newcommand{\recorda}[1]{\par\smallskip\begin{tcolorbox}[colback=accent!7,colframe=accent!55,arc=2mm,boxrule=0.5pt,left=2.5mm,right=2.5mm,top=1mm,bottom=1mm]\small{\color{accent}\faLightbulb[regular]}\ \textbf{Recuerda:}\ #1\end{tcolorbox}}
\newcommand{\res}[1]{\boxed{#1}}
\newcommand{\rest}[1]{\fbox{#1}}

\def\fmneed{14}
% =====================================================================
% Capçalera FentMates de les fitxes: un color per a cada curs.
% Els colors són els de la pàgina inicial (LEVELS a assets/js/materials.js);
% el logo ve de _eines/logo/pdf/ (python3 _eines/logo/genera.py).
%
% Ús, al final del preàmbul de la unitat:
%   % =====================================================================
% Capçalera FentMates de les fitxes: un color per a cada curs.
% Els colors són els de la pàgina inicial (LEVELS a assets/js/materials.js);
% el logo ve de _eines/logo/pdf/ (python3 _eines/logo/genera.py).
%
% Ús, al final del preàmbul de la unitat:
%   % =====================================================================
% Capçalera FentMates de les fitxes: un color per a cada curs.
% Els colors són els de la pàgina inicial (LEVELS a assets/js/materials.js);
% el logo ve de _eines/logo/pdf/ (python3 _eines/logo/genera.py).
%
% Ús, al final del preàmbul de la unitat:
%   \input{../comu/fentmates.tex}
%   \fmcurs{1eso}
% i, al començament del document:
%   \capcalera{Matemàtiques}{1r ESO}{Fitxa d'exercicis · Unitat 2}{Divisibilitat}
% Opcional abans de \input: \def\fmneed{14} (espai mínim abans de cada secció).
% =====================================================================
\makeatletter
\def\fm@logodir{../../logo/pdf}
\providecommand{\fmneed}{10}
\def\fm@cursnom{}

% clau del curs, color i color del text damunt del color
\newcommand{\fm@curs}[3]{\@namedef{fm@color@#1}{#2}\@namedef{fm@text@#1}{#3}}
\fm@curs{1eso}{E6B84A}{navy}
\fm@curs{2eso}{3B82F6}{white}
\fm@curs{3eso}{FF7A50}{white}
\fm@curs{4eso-a}{A78BFA}{white}
\fm@curs{4eso-b}{6D28D9}{white}
\fm@curs{1batx-mi}{9B1C1C}{white}
\fm@curs{1batx-ccss}{EF6B6B}{white}
\fm@curs{2batx-mii}{1E6B47}{white}
\fm@curs{2batx-ccss}{2BB3A3}{white}

\newcommand{\fmcurs}[1]{%
  \@ifundefined{fm@color@#1}{\PackageError{fentmates}{Curs desconegut: #1}{}}{}%
  \edef\fm@tmp{\noexpand\definecolor{curs}{HTML}{\@nameuse{fm@color@#1}}%
               \noexpand\colorlet{curstext}{\@nameuse{fm@text@#1}}}\fm@tmp
  \colorlet{cursfosc}{curs!72!black}%
  \def\fm@logo{\fm@logodir/vertical-#1.pdf}%
  \def\fm@marca{\fm@logodir/marca-#1.pdf}}

% ---------- capçalera de la primera pàgina ----------
% #1 matèria (petit), #2 curs, #3 tipus de document i unitat, #4 títol
\newcommand{\capcalera}[4]{%
  \gdef\fm@cursnom{#2}%
  \par\noindent
  \begin{tcolorbox}[colback=curs,colframe=curs,boxrule=0pt,arc=2.2mm,width=3.5cm,height=2.15cm,
      valign=center,left=3mm,right=2mm,top=1mm,bottom=1mm,nobeforeafter,box align=center]
    \color{curstext}{\footnotesize #1}\par\vspace{1pt}{\fontsize{19}{21}\selectfont\bfseries #2\par}%
  \end{tcolorbox}\hspace{5mm}%
  \begin{minipage}[c]{\dimexpr\linewidth-3.5cm-5mm-2.9cm\relax}
    {\color{cursfosc}\small\bfseries\MakeUppercase{#3}}\par\vspace{2pt}
    {\color{navy}\fontsize{20}{23}\selectfont\bfseries #4\par}
  \end{minipage}\hfill
  \begin{minipage}[c]{2.6cm}\raggedleft\includegraphics[width=2.6cm]{\fm@logo}\end{minipage}\par
  \vspace{2pt}}

% ---------- barres de secció del color del curs ----------
\renewcommand{\seccio}[2]{\par\needspace{\fmneed\baselineskip}\bigskip
  \begin{tcolorbox}[colback=curs,colframe=curs,coltext=curstext,arc=2mm,boxrule=0pt,left=3mm,right=3mm,top=1.2mm,bottom=1.2mm]
  {\large\bfseries #1\quad #2}
  \end{tcolorbox}\vspace{-2pt}}

% ---------- estrelles de dificultat del color del curs (a Batxillerat no se'n posen) ----------
\@ifundefined{estrelles}{}{%
  \renewcommand{\estrelles}[1]{{\color{curs}\ifcase#1\or\faStar\faStar[regular]\faStar[regular]\or\faStar\faStar\faStar[regular]\or\faStar\faStar\faStar\fi}}}

% ---------- capçalera i peu de les pàgines ----------
% primera pàgina: logos institucionals al peu
\fancypagestyle{primera}{\fancyhf{}\renewcommand{\headrulewidth}{0pt}
  \fancyfoot[L]{\raisebox{-0.32\height}{\includegraphics[height=0.78cm]{logo-gva.png}}\hspace{5mm}%
    \raisebox{-0.32\height}{\includegraphics[height=0.78cm]{logo-vermellar.png}}}
  \fancyfoot[R]{\scriptsize\color{muted}\href{https://creativecommons.org/licenses/by-nc-sa/4.0/deed.ca}{\ccbyncsa}\ \ CC BY-NC-SA 4.0 · fentmates.org\qquad\small\thepage}}
% la resta: el símbol a l'esquerra i el curs a la dreta, amb el títol al centre com abans
\fancyhead[L]{\raisebox{-0.3\height}{\includegraphics[height=0.62cm]{\fm@marca}}}
\fancyhead[R]{\small\bfseries\color{cursfosc}\fm@cursnom}
\renewcommand{\headrule}{\color{curs!45}\hrule height \headrulewidth}
\makeatother

%   \fmcurs{1eso}
% i, al començament del document:
%   \capcalera{Matemàtiques}{1r ESO}{Fitxa d'exercicis · Unitat 2}{Divisibilitat}
% Opcional abans de \input: \def\fmneed{14} (espai mínim abans de cada secció).
% =====================================================================
\makeatletter
\def\fm@logodir{../../logo/pdf}
\providecommand{\fmneed}{10}
\def\fm@cursnom{}

% clau del curs, color i color del text damunt del color
\newcommand{\fm@curs}[3]{\@namedef{fm@color@#1}{#2}\@namedef{fm@text@#1}{#3}}
\fm@curs{1eso}{E6B84A}{navy}
\fm@curs{2eso}{3B82F6}{white}
\fm@curs{3eso}{FF7A50}{white}
\fm@curs{4eso-a}{A78BFA}{white}
\fm@curs{4eso-b}{6D28D9}{white}
\fm@curs{1batx-mi}{9B1C1C}{white}
\fm@curs{1batx-ccss}{EF6B6B}{white}
\fm@curs{2batx-mii}{1E6B47}{white}
\fm@curs{2batx-ccss}{2BB3A3}{white}

\newcommand{\fmcurs}[1]{%
  \@ifundefined{fm@color@#1}{\PackageError{fentmates}{Curs desconegut: #1}{}}{}%
  \edef\fm@tmp{\noexpand\definecolor{curs}{HTML}{\@nameuse{fm@color@#1}}%
               \noexpand\colorlet{curstext}{\@nameuse{fm@text@#1}}}\fm@tmp
  \colorlet{cursfosc}{curs!72!black}%
  \def\fm@logo{\fm@logodir/vertical-#1.pdf}%
  \def\fm@marca{\fm@logodir/marca-#1.pdf}}

% ---------- capçalera de la primera pàgina ----------
% #1 matèria (petit), #2 curs, #3 tipus de document i unitat, #4 títol
\newcommand{\capcalera}[4]{%
  \gdef\fm@cursnom{#2}%
  \par\noindent
  \begin{tcolorbox}[colback=curs,colframe=curs,boxrule=0pt,arc=2.2mm,width=3.5cm,height=2.15cm,
      valign=center,left=3mm,right=2mm,top=1mm,bottom=1mm,nobeforeafter,box align=center]
    \color{curstext}{\footnotesize #1}\par\vspace{1pt}{\fontsize{19}{21}\selectfont\bfseries #2\par}%
  \end{tcolorbox}\hspace{5mm}%
  \begin{minipage}[c]{\dimexpr\linewidth-3.5cm-5mm-2.9cm\relax}
    {\color{cursfosc}\small\bfseries\MakeUppercase{#3}}\par\vspace{2pt}
    {\color{navy}\fontsize{20}{23}\selectfont\bfseries #4\par}
  \end{minipage}\hfill
  \begin{minipage}[c]{2.6cm}\raggedleft\includegraphics[width=2.6cm]{\fm@logo}\end{minipage}\par
  \vspace{2pt}}

% ---------- barres de secció del color del curs ----------
\renewcommand{\seccio}[2]{\par\needspace{\fmneed\baselineskip}\bigskip
  \begin{tcolorbox}[colback=curs,colframe=curs,coltext=curstext,arc=2mm,boxrule=0pt,left=3mm,right=3mm,top=1.2mm,bottom=1.2mm]
  {\large\bfseries #1\quad #2}
  \end{tcolorbox}\vspace{-2pt}}

% ---------- estrelles de dificultat del color del curs (a Batxillerat no se'n posen) ----------
\@ifundefined{estrelles}{}{%
  \renewcommand{\estrelles}[1]{{\color{curs}\ifcase#1\or\faStar\faStar[regular]\faStar[regular]\or\faStar\faStar\faStar[regular]\or\faStar\faStar\faStar\fi}}}

% ---------- capçalera i peu de les pàgines ----------
% primera pàgina: logos institucionals al peu
\fancypagestyle{primera}{\fancyhf{}\renewcommand{\headrulewidth}{0pt}
  \fancyfoot[L]{\raisebox{-0.32\height}{\includegraphics[height=0.78cm]{logo-gva.png}}\hspace{5mm}%
    \raisebox{-0.32\height}{\includegraphics[height=0.78cm]{logo-vermellar.png}}}
  \fancyfoot[R]{\scriptsize\color{muted}\href{https://creativecommons.org/licenses/by-nc-sa/4.0/deed.ca}{\ccbyncsa}\ \ CC BY-NC-SA 4.0 · fentmates.org\qquad\small\thepage}}
% la resta: el símbol a l'esquerra i el curs a la dreta, amb el títol al centre com abans
\fancyhead[L]{\raisebox{-0.3\height}{\includegraphics[height=0.62cm]{\fm@marca}}}
\fancyhead[R]{\small\bfseries\color{cursfosc}\fm@cursnom}
\renewcommand{\headrule}{\color{curs!45}\hrule height \headrulewidth}
\makeatother

%   \fmcurs{1eso}
% i, al començament del document:
%   \capcalera{Matemàtiques}{1r ESO}{Fitxa d'exercicis · Unitat 2}{Divisibilitat}
% Opcional abans de \input: \def\fmneed{14} (espai mínim abans de cada secció).
% =====================================================================
\makeatletter
\def\fm@logodir{../../logo/pdf}
\providecommand{\fmneed}{10}
\def\fm@cursnom{}

% clau del curs, color i color del text damunt del color
\newcommand{\fm@curs}[3]{\@namedef{fm@color@#1}{#2}\@namedef{fm@text@#1}{#3}}
\fm@curs{1eso}{E6B84A}{navy}
\fm@curs{2eso}{3B82F6}{white}
\fm@curs{3eso}{FF7A50}{white}
\fm@curs{4eso-a}{A78BFA}{white}
\fm@curs{4eso-b}{6D28D9}{white}
\fm@curs{1batx-mi}{9B1C1C}{white}
\fm@curs{1batx-ccss}{EF6B6B}{white}
\fm@curs{2batx-mii}{1E6B47}{white}
\fm@curs{2batx-ccss}{2BB3A3}{white}

\newcommand{\fmcurs}[1]{%
  \@ifundefined{fm@color@#1}{\PackageError{fentmates}{Curs desconegut: #1}{}}{}%
  \edef\fm@tmp{\noexpand\definecolor{curs}{HTML}{\@nameuse{fm@color@#1}}%
               \noexpand\colorlet{curstext}{\@nameuse{fm@text@#1}}}\fm@tmp
  \colorlet{cursfosc}{curs!72!black}%
  \def\fm@logo{\fm@logodir/vertical-#1.pdf}%
  \def\fm@marca{\fm@logodir/marca-#1.pdf}}

% ---------- capçalera de la primera pàgina ----------
% #1 matèria (petit), #2 curs, #3 tipus de document i unitat, #4 títol
\newcommand{\capcalera}[4]{%
  \gdef\fm@cursnom{#2}%
  \par\noindent
  \begin{tcolorbox}[colback=curs,colframe=curs,boxrule=0pt,arc=2.2mm,width=3.5cm,height=2.15cm,
      valign=center,left=3mm,right=2mm,top=1mm,bottom=1mm,nobeforeafter,box align=center]
    \color{curstext}{\footnotesize #1}\par\vspace{1pt}{\fontsize{19}{21}\selectfont\bfseries #2\par}%
  \end{tcolorbox}\hspace{5mm}%
  \begin{minipage}[c]{\dimexpr\linewidth-3.5cm-5mm-2.9cm\relax}
    {\color{cursfosc}\small\bfseries\MakeUppercase{#3}}\par\vspace{2pt}
    {\color{navy}\fontsize{20}{23}\selectfont\bfseries #4\par}
  \end{minipage}\hfill
  \begin{minipage}[c]{2.6cm}\raggedleft\includegraphics[width=2.6cm]{\fm@logo}\end{minipage}\par
  \vspace{2pt}}

% ---------- barres de secció del color del curs ----------
\renewcommand{\seccio}[2]{\par\needspace{\fmneed\baselineskip}\bigskip
  \begin{tcolorbox}[colback=curs,colframe=curs,coltext=curstext,arc=2mm,boxrule=0pt,left=3mm,right=3mm,top=1.2mm,bottom=1.2mm]
  {\large\bfseries #1\quad #2}
  \end{tcolorbox}\vspace{-2pt}}

% ---------- estrelles de dificultat del color del curs (a Batxillerat no se'n posen) ----------
\@ifundefined{estrelles}{}{%
  \renewcommand{\estrelles}[1]{{\color{curs}\ifcase#1\or\faStar\faStar[regular]\faStar[regular]\or\faStar\faStar\faStar[regular]\or\faStar\faStar\faStar\fi}}}

% ---------- capçalera i peu de les pàgines ----------
% primera pàgina: logos institucionals al peu
\fancypagestyle{primera}{\fancyhf{}\renewcommand{\headrulewidth}{0pt}
  \fancyfoot[L]{\raisebox{-0.32\height}{\includegraphics[height=0.78cm]{logo-gva.png}}\hspace{5mm}%
    \raisebox{-0.32\height}{\includegraphics[height=0.78cm]{logo-vermellar.png}}}
  \fancyfoot[R]{\scriptsize\color{muted}\href{https://creativecommons.org/licenses/by-nc-sa/4.0/deed.ca}{\ccbyncsa}\ \ CC BY-NC-SA 4.0 · fentmates.org\qquad\small\thepage}}
% la resta: el símbol a l'esquerra i el curs a la dreta, amb el títol al centre com abans
\fancyhead[L]{\raisebox{-0.3\height}{\includegraphics[height=0.62cm]{\fm@marca}}}
\fancyhead[R]{\small\bfseries\color{cursfosc}\fm@cursnom}
\renewcommand{\headrule}{\color{curs!45}\hrule height \headrulewidth}
\makeatother

\fmcurs{3eso}

\begin{document}
\thispagestyle{primera}

\capcalera{Matemáticas}{3\textsuperscript{o} ESO}{Soluciones de la ficha · Unidad 12}{Probabilidad}
\vspace{8pt}

\begin{tcolorbox}[colback=soft,colframe=soft,arc=2mm,left=3mm,right=3mm,top=1.5mm,bottom=1.5mm]
\small Intenta resolver cada ejercicio antes de mirar su solución. Si no te sale el mismo resultado, repasa los pasos o consulta los apuntes interactivos de la unidad: encontrarás simulaciones, una bolsa de bolas, la tabla de los dos dados y un diagrama de árbol que calcula las probabilidades paso a paso. Las probabilidades se dan en fracción irreducible (o en decimal, si los datos lo son).
\end{tcolorbox}

% =====================================================
\seccio{1}{Experimentos aleatorios. Espacio muestral y sucesos}

\ex{1} \begin{tasks}(2)
  \task \rest{aleatorio}
  \task \rest{determinista}
  \task \rest{aleatorio}
  \task \rest{determinista} (6 €)
  \task \rest{aleatorio}
  \task \rest{determinista}
\end{tasks}

\ex{1} \begin{tasks}(2)
  \task $E=\{C,\,X\}$
  \task $E=\{1,\,2,\,3,\,4\}$
  \task $E=\{\text{rojo},\,\text{azul},\,\text{verde},\,\text{amarillo},\,\text{negro}\}$
  \task $E=\{R,\,A\}$
  \task $E=\{\text{lu},\,\text{ma},\,\text{mi},\,\text{ju},\,\text{vi},\,\text{sá},\,\text{do}\}$
  \task $E=\{0,\,1,\,2,\,3\}$
\end{tasks}

\ex{1} \begin{tasks}(3)
  \task $\{2,4,6,8,10,12\}$: \rest{compuesto}
  \task $\{5,\,10\}$: \rest{compuesto}
  \task $\varnothing$: \rest{imposible}
  \task $\{1,2,3,4,6,12\}$: \rest{compuesto}
  \task $E$: \rest{seguro}
  \task $\{11\}$: \rest{elemental}
\end{tasks}

\ex{2} \par\smallskip\begin{center}\begin{tikzpicture}[grow'=right,level distance=14mm,level 1/.style={sibling distance=13mm},level 2/.style={sibling distance=6.5mm},every node/.style={font=\small\bfseries},edge from parent/.style={draw=muted,thick}]
\node[circle,fill=navy,inner sep=1.2pt] {} child {node[] {C} child {node[label={[font=\scriptsize,xshift=1mm]right:$CC$}] {C} edge from parent node[above,pos=.55,font=\scriptsize,text=teal!70!black,inner sep=1.5pt] {}} child {node[label={[font=\scriptsize,xshift=1mm]right:$CX$}] {X} edge from parent node[below,pos=.55,font=\scriptsize,text=teal!70!black,inner sep=1.5pt] {}} edge from parent node[above,pos=.55,font=\scriptsize,text=teal!70!black,inner sep=1.5pt] {}} child {node[] {X} child {node[label={[font=\scriptsize,xshift=1mm]right:$XC$}] {C} edge from parent node[above,pos=.55,font=\scriptsize,text=teal!70!black,inner sep=1.5pt] {}} child {node[label={[font=\scriptsize,xshift=1mm]right:$XX$}] {X} edge from parent node[below,pos=.55,font=\scriptsize,text=teal!70!black,inner sep=1.5pt] {}} edge from parent node[below,pos=.55,font=\scriptsize,text=teal!70!black,inner sep=1.5pt] {}};
\end{tikzpicture}\end{center}$E=\{CC,\,CX,\,XC,\,XX\}$; \quad $A=\res{\{CC,\,CX,\,XC\}}$, \quad $B=\res{\{XX\}}$, \quad $C=\res{\{CC,\,XX\}}$.

\ex{2} \par\smallskip\begin{center}{\small\renewcommand{\arraystretch}{1.2}\begin{tabular}{l|cccccc}
\rowcolor{soft}  & 1 & 2 & 3 & 4 & 5 & 6\\\hline
$C$ & $C1$ & $C2$ & $C3$ & $C4$ & $C5$ & $C6$\\
$X$ & $X1$ & $X2$ & $X3$ & $X4$ & $X5$ & $X6$\\
\end{tabular}}\end{center}$2\cdot6=\res{12}$ resultados. \quad «Cara y par» $=\res{\{C2,\,C4,\,C6\}}$; \quad «cruz y menor que 3» $=\res{\{X1,\,X2\}}$.

\ex{2} \par\smallskip\begin{center}\begin{tikzpicture}[grow'=right,level distance=13mm,level 1/.style={sibling distance=20mm},level 2/.style={sibling distance=10mm},level 3/.style={sibling distance=5mm},every node/.style={font=\small\bfseries},edge from parent/.style={draw=muted,thick}]
\node[circle,fill=navy,inner sep=1.2pt] {} child {node[] {O} child {node[] {O} child {node[label={[font=\scriptsize,xshift=1mm]right:$OOO$}] {O} edge from parent node[above,pos=.55,font=\scriptsize,text=teal!70!black,inner sep=1.5pt] {}} child {node[label={[font=\scriptsize,xshift=1mm]right:$OOA$}] {A} edge from parent node[below,pos=.55,font=\scriptsize,text=teal!70!black,inner sep=1.5pt] {}} edge from parent node[above,pos=.55,font=\scriptsize,text=teal!70!black,inner sep=1.5pt] {}} child {node[] {A} child {node[label={[font=\scriptsize,xshift=1mm]right:$OAO$}] {O} edge from parent node[above,pos=.55,font=\scriptsize,text=teal!70!black,inner sep=1.5pt] {}} child {node[label={[font=\scriptsize,xshift=1mm]right:$OAA$}] {A} edge from parent node[below,pos=.55,font=\scriptsize,text=teal!70!black,inner sep=1.5pt] {}} edge from parent node[below,pos=.55,font=\scriptsize,text=teal!70!black,inner sep=1.5pt] {}} edge from parent node[above,pos=.55,font=\scriptsize,text=teal!70!black,inner sep=1.5pt] {}} child {node[] {A} child {node[] {O} child {node[label={[font=\scriptsize,xshift=1mm]right:$AOO$}] {O} edge from parent node[above,pos=.55,font=\scriptsize,text=teal!70!black,inner sep=1.5pt] {}} child {node[label={[font=\scriptsize,xshift=1mm]right:$AOA$}] {A} edge from parent node[below,pos=.55,font=\scriptsize,text=teal!70!black,inner sep=1.5pt] {}} edge from parent node[above,pos=.55,font=\scriptsize,text=teal!70!black,inner sep=1.5pt] {}} child {node[] {A} child {node[label={[font=\scriptsize,xshift=1mm]right:$AAO$}] {O} edge from parent node[above,pos=.55,font=\scriptsize,text=teal!70!black,inner sep=1.5pt] {}} child {node[label={[font=\scriptsize,xshift=1mm]right:$AAA$}] {A} edge from parent node[below,pos=.55,font=\scriptsize,text=teal!70!black,inner sep=1.5pt] {}} edge from parent node[below,pos=.55,font=\scriptsize,text=teal!70!black,inner sep=1.5pt] {}} edge from parent node[below,pos=.55,font=\scriptsize,text=teal!70!black,inner sep=1.5pt] {}};
\end{tikzpicture}\end{center}$E=\{OOO,\,OOA,\,OAO,\,OAA,\,AOO,\,AOA,\,AAO,\,AAA\}$ (8 resultados). \quad «Dos chicas» $=\res{\{OAA,\,AOA,\,AAO\}}$.

\ex{3} $E=\{123,\,132,\,213,\,231,\,312,\,321\}$. \quad «Par» $=\res{\{132,\,312\}}$; \quad «mayor que 200» $=\res{\{213,\,231,\,312,\,321\}}$.

\ex{3} $6\cdot6=\res{36}$ resultados (pares ordenados). «Suma 5» $=\res{\{(1,4),\,(2,3),\,(3,2),\,(4,1)\}}$; «doble» $=\res{\{(1,1),\,(2,2),\,\dots,\,(6,6)\}}$ (6 resultados).


% =====================================================
\seccio{2}{Probabilidad de un suceso. Propiedades}

\ex{1} $h$: \ $\res{0{,}14}$, \ $\res{0{,}18}$, \ $\res{0{,}16}$, \ $\res{0{,}2}$, \ $\res{0{,}12}$, \ $\res{0{,}2}$ \ (suman 1). \quad $h(\text{par})=\frac{9+10+10}{50}=\res{0{,}58}$.

\ex{1} Una probabilidad siempre está entre 0 y 1. Pueden serlo: \rest{$0{,}7$, \ $0$, \ $100\,\%$ y $\frac{99}{100}$}. No pueden serlo: $1{,}2$, $-0{,}1$, $\frac32$ y $\frac54$.

\ex{1} $P(\text{verde})=1-(0{,}5+0{,}3)=\res{0{,}2}$. \ $P(\text{rojo o verde})=0{,}5+0{,}2=\res{0{,}7}$.

\ex{2} $P(\text{amarillo})=1-(0{,}35+0{,}25+0{,}1)=\res{0{,}3}$. \ $P(\text{no rojo})=P(\text{azul})+P(\text{verde})+P(\text{amarillo})=0{,}25+0{,}1+0{,}3=\res{0{,}65}$. \ $P(\text{azul o amarillo})=0{,}25+0{,}3=\res{0{,}55}$.

\ex{2} Si $P(1)=\dots=P(5)=x$: \ $5x+0{,}5=1\Rightarrow x=\res{0{,}1}$. \ $P(\text{par})=0{,}1+0{,}1+0{,}5=\res{0{,}7}$. \ $P(\text{menos de 3})=0{,}1+0{,}1=\res{0{,}2}$.

\ex{3} Las seis probabilidades suman 1: $p+2p+\dots+6p=21p=1\Rightarrow p=\res{\frac1{21}}$. \ $P(6)=\frac6{21}=\res{\frac27}$. \ $P(\text{par})=\frac{2+4+6}{21}=\frac{12}{21}=\res{\frac47}$.

\ex{3} Si $P(\text{azul})=x$: \ $2x+x+0{,}4=1\Rightarrow x=0{,}2$. \ $P(\text{roja})=\res{0{,}4}$, \ $P(\text{azul})=\res{0{,}2}$. \ Con 20 bolas: \res{8 rojas, 4 azules y 8 verdes}.


% =====================================================
\seccio{3}{Experimentos regulares. Ley de Laplace}

\ex{1} \begin{tasks}(3)
  \task $\res{\frac16}$
  \task $\frac36=\res{\frac12}$
  \task $\frac46=\res{\frac23}$
  \task $\frac26=\res{\frac13}$
  \task $\res{0}$ (imposible)
  \task $\res{1}$ (seguro)
\end{tasks}

\ex{1} Hay 15 bolas: \ $P(R)=\frac6{15}=\res{\frac25}$, \ $P(A)=\res{\frac4{15}}$, \ $P(V)=\frac5{15}=\res{\frac13}$, \ $P(\text{no verde})=\frac{6+4}{15}=\frac{10}{15}=\res{\frac23}$.

\ex{1} a) $\{4,8,12\}$: $\frac3{12}=\res{\frac14}$. \ b) $\{2,3,5,7,11\}$: $\res{\frac5{12}}$. \ c) $\{1,2,3,4,6,12\}$: $\frac6{12}=\res{\frac12}$.

\ex{2} \begin{tasks}(3)
  \task $\frac{10}{40}=\res{\frac14}$
  \task $\frac4{40}=\res{\frac1{10}}$
  \task $\res{\frac3{40}}$
  \task $\frac{28}{40}=\res{\frac7{10}}$
  \task $\frac{10+4-1}{40}=\res{\frac{13}{40}}$
  \task $\res{\frac1{40}}$
\end{tasks}

\ex{2} 11 letras; vocales A, E, A, I, A: \ $P(\text{vocal})=\res{\frac5{11}}$; \ $P(M)=\res{\frac2{11}}$; \ $P(\text{consonante})=\frac{11-5}{11}=\res{\frac6{11}}$.

\ex{2} a) $07,17,\dots,97$: $\frac{10}{100}=\res{\frac1{10}}$. \ b) $00,11,\dots,99$: $\res{\frac1{10}}$. \ c) $90,\dots,99$: $\res{\frac1{10}}$.

\ex{2} $\{3,5,6,9,10,12,15,18,20\}$: 9 casos ($6+4-1$, porque el 15 es múltiplo de los dos): \ $P=\res{\frac9{20}}$.

\ex{2} $P(\text{Pablo})=\frac26=\frac13$ y $P(\text{María})=\frac46=\frac23$: \rest{no es justo}. Por ejemplo: Pablo gana con par y María con impar ($\frac12$ cada uno).

\ex{3} Si hay $x$ azules: $\frac{x}{5+x}=\frac23\Rightarrow3x=10+2x\Rightarrow\res{x=10}$. \ (Comprobación: $\frac{10}{15}=\frac23$ \checkmark)

\ex{3} $P(V)=1-\frac14-\frac13=\frac5{12}$; si hay $n$ bolas, $\frac{10}{n}=\frac5{12}\Rightarrow n=24$: \ \res{6 rojas, 8 azules y 10 verdes}.

\ex{3} \rest{No} es regular: los sectores no son iguales. La probabilidad es la fracción de círculo que ocupa cada sector: \ $P(\text{rojo})=\frac{180}{360}=\res{\frac12}$, \ $P(\text{azul})=\frac{120}{360}=\res{\frac13}$, \ $P(\text{verde})=\frac{60}{360}=\res{\frac16}$.


% =====================================================
\seccio{4}{Experimentos irregulares. Ley de los grandes números}

\ex{1} \begin{tasks}(2)
  \task \rest{regular}
  \task \rest{irregular}
  \task \rest{irregular}
  \task \rest{regular}
  \task \rest{irregular}
  \task \rest{irregular}
\end{tasks}

\ex{1} $h$: \ $0{,}65$; \ $0{,}58$; \ $0{,}594$; \ $0{,}601$. \ Con más lanzamientos se estabiliza: $P(\text{punta arriba})\approx\res{0{,}6}$.

\ex{2} $P\approx\frac{45}{60}=\res{0{,}75}$. \ $20\cdot0{,}75=\res{15\text{ penaltis}}$ (aproximadamente).

\ex{2} \rest{No} necesariamente: 10 lanzamientos son muy pocos y la frecuencia relativa varía mucho. Para saber si está trucada habría que lanzarla muchas más veces (cientos o miles) y ver si $h$ se acerca a 0,5.

\ex{2} $600\cdot\frac16=\res{100}$ veces cada uno. Los resultados están todos cerca de 100: \rest{no parece trucado}.

\ex{2} $P(R)\approx\frac{120}{300}=0{,}4$; \ si hay $x$ rojas, $\frac{x}{30}\approx0{,}4\Rightarrow x\approx\res{12\text{ rojas}}$ (y 18 blancas).

\ex{3} $P(6)\approx\frac{400}{1200}=\res{\frac13}$; \ $P(\text{par})\approx\frac{150+150+400}{1200}=\frac{700}{1200}=\res{\frac7{12}}\approx0{,}58$. \ $300\cdot\frac13=\res{100\text{ seises}}$. \ \rest{Sí}: en un dado correcto esperaríamos unos 200 de cada resultado.


% =====================================================
\seccio{5}{Experimentos compuestos. Diagramas de árbol}

\ex{1} a) $\frac12\cdot\frac16=\res{\frac1{12}}$. \ b) $\frac12\cdot\frac36=\res{\frac14}$. \ c) $\{2,3,5\}$: $\frac36=\res{\frac12}$.

\ex{1} \par\smallskip\begin{center}\begin{tikzpicture}[grow'=right,level distance=16mm,level 1/.style={sibling distance=14mm},level 2/.style={sibling distance=7mm},every node/.style={font=\small\bfseries},edge from parent/.style={draw=muted,thick}]
\node[circle,fill=navy,inner sep=1.2pt] {} child {node[] {C} child {node[label={[font=\scriptsize,xshift=1mm]right:$CC$}] {C} edge from parent node[above,pos=.55,font=\scriptsize,text=teal!70!black,inner sep=1.5pt] {$\frac12$}} child {node[label={[font=\scriptsize,xshift=1mm]right:$CX$}] {X} edge from parent node[below,pos=.55,font=\scriptsize,text=teal!70!black,inner sep=1.5pt] {$\frac12$}} edge from parent node[above,pos=.55,font=\scriptsize,text=teal!70!black,inner sep=1.5pt] {$\frac12$}} child {node[] {X} child {node[label={[font=\scriptsize,xshift=1mm]right:$XC$}] {C} edge from parent node[above,pos=.55,font=\scriptsize,text=teal!70!black,inner sep=1.5pt] {$\frac12$}} child {node[label={[font=\scriptsize,xshift=1mm]right:$XX$}] {X} edge from parent node[below,pos=.55,font=\scriptsize,text=teal!70!black,inner sep=1.5pt] {$\frac12$}} edge from parent node[below,pos=.55,font=\scriptsize,text=teal!70!black,inner sep=1.5pt] {$\frac12$}};
\end{tikzpicture}\end{center}$P(CC)=\frac12\cdot\frac12=\res{\frac14}$. \ $P(\text{al menos una cruz})=P(CX)+P(XC)+P(XX)=\frac14+\frac14+\frac14=\res{\frac34}$.

\ex{1} \par\smallskip\begin{center}{\small\renewcommand{\arraystretch}{1.15}\setlength{\tabcolsep}{4pt}\begin{tabular}{c|cccccc}
\rowcolor{soft} $+$ & 1 & 2 & 3 & 4 & 5 & 6\\\hline
\cellcolor{soft}1 & 2 & 3 & 4 & 5 & 6 & 7\\
\cellcolor{soft}2 & 3 & 4 & 5 & 6 & 7 & \cellcolor{accent!30}\textbf{8}\\
\cellcolor{soft}3 & 4 & 5 & 6 & 7 & \cellcolor{accent!30}\textbf{8} & 9\\
\cellcolor{soft}4 & 5 & 6 & 7 & \cellcolor{accent!30}\textbf{8} & 9 & 10\\
\cellcolor{soft}5 & 6 & 7 & \cellcolor{accent!30}\textbf{8} & 9 & 10 & 11\\
\cellcolor{soft}6 & 7 & \cellcolor{accent!30}\textbf{8} & 9 & 10 & 11 & 12\\
\end{tabular}}\end{center}a) $\res{\frac5{36}}$ (marcadas). \ b) $\res{\frac1{36}}$. \ c) $\frac{3+2+1}{36}=\res{\frac16}$. \ d) $\frac{18}{36}=\res{\frac12}$.

\ex{2} a) $\frac6{36}=\res{\frac16}$. \ b) $(1,3),(2,4),(3,5),(4,6)$ y al revés: $\frac8{36}=\res{\frac29}$. \ c) En la tabla de 36 productos, 27 son pares (solo son impares los 9 de dos números impares): $\frac{27}{36}=\res{\frac34}$.

\ex{2} Cada camino del árbol tiene probabilidad $\frac12\cdot\frac12\cdot\frac12=\frac18$. \ a) $\res{\frac18}$. \ b) $CXX$, $XCX$, $XXC$: $\res{\frac38}$. \ c) Todos los caminos menos $XXX$: $\res{\frac78}$.

\ex{2} \par\smallskip\begin{center}\begin{tikzpicture}[grow'=right,level distance=17mm,level 1/.style={sibling distance=15mm},level 2/.style={sibling distance=7.5mm},every node/.style={font=\small\bfseries},edge from parent/.style={draw=muted,thick}]
\node[circle,fill=navy,inner sep=1.2pt] {} child {node[] {G} child {node[label={[font=\scriptsize,xshift=1mm]right:$GG$: $0{,}64$}] {G} edge from parent node[above,pos=.55,font=\scriptsize,text=teal!70!black,inner sep=1.5pt] {$0{,}8$}} child {node[label={[font=\scriptsize,xshift=1mm]right:$GF$: $0{,}16$}] {F} edge from parent node[below,pos=.55,font=\scriptsize,text=teal!70!black,inner sep=1.5pt] {$0{,}2$}} edge from parent node[above,pos=.55,font=\scriptsize,text=teal!70!black,inner sep=1.5pt] {$0{,}8$}} child {node[] {F} child {node[label={[font=\scriptsize,xshift=1mm]right:$FG$: $0{,}16$}] {G} edge from parent node[above,pos=.55,font=\scriptsize,text=teal!70!black,inner sep=1.5pt] {$0{,}8$}} child {node[label={[font=\scriptsize,xshift=1mm]right:$FF$: $0{,}04$}] {F} edge from parent node[below,pos=.55,font=\scriptsize,text=teal!70!black,inner sep=1.5pt] {$0{,}2$}} edge from parent node[below,pos=.55,font=\scriptsize,text=teal!70!black,inner sep=1.5pt] {$0{,}2$}};
\end{tikzpicture}\end{center}$P(GG)=\res{0{,}64}$; \ $P(FF)=\res{0{,}04}$; \ $P(\text{exactamente uno})=0{,}16+0{,}16=\res{0{,}32}$. \ (Suman 1 \checkmark)

\ex{2} a) $\frac16\cdot\frac16=\res{\frac1{36}}$. \ b) $\frac56\cdot\frac56=\res{\frac{25}{36}}$. \ c) En la tabla, la fila y la columna del 6 tienen 11 casillas: $\res{\frac{11}{36}}$.

\ex{3} $P(\text{todas})=\left(\frac14\right)^3=\res{\frac1{64}}$. \ $P(\text{ninguna})=\left(\frac34\right)^3=\res{\frac{27}{64}}$. \ $P(\text{exactamente una})=3\cdot\frac14\cdot\frac34\cdot\frac34=\res{\frac{27}{64}}$ (tres caminos: el acierto puede estar en la primera, la segunda o la tercera pregunta).

\ex{3} El primero puede ser cualquiera; los otros dos tienen que coincidir: $1\cdot\frac16\cdot\frac16=\res{\frac1{36}}$. \ $P(\text{ningún 6})=\frac56\cdot\frac56\cdot\frac56=\res{\frac{125}{216}}\approx0{,}58$.

\ex{3} Contamos en la tabla: a) $\frac6{36}=\res{\frac16}$. \ b) $\res{\frac{11}{36}}$. \ c) $(1,6)$ y $(6,1)$: $\frac2{36}=\res{\frac1{18}}$. \ d) $6+11-2=15$ casillas (las 2 de c) no se cuentan dos veces): $\frac{15}{36}=\res{\frac5{12}}$.


% =====================================================
\seccio{6}{Extracciones con y sin reemplazamiento. Problemas}

\ex{1} \par\smallskip\begin{center}\begin{tikzpicture}[grow'=right,level distance=17mm,level 1/.style={sibling distance=16mm},level 2/.style={sibling distance=8mm},every node/.style={font=\small\bfseries},edge from parent/.style={draw=muted,thick}]
\node[circle,fill=navy,inner sep=1.2pt] {} child {node[circle,fill=blue!70!black,text=white,inner sep=1.2pt,minimum size=4.2mm,font=\scriptsize\bfseries] {B} child {node[circle,fill=blue!70!black,text=white,inner sep=1.2pt,minimum size=4.2mm,font=\scriptsize\bfseries,label={[font=\scriptsize,xshift=1mm]right:$BB$: $\frac{9}{25}$}] {B} edge from parent node[above,pos=.55,font=\scriptsize,text=teal!70!black,inner sep=1.5pt] {$\frac{3}{5}$}} child {node[circle,fill=black!80,text=white,inner sep=1.2pt,minimum size=4.2mm,font=\scriptsize\bfseries,label={[font=\scriptsize,xshift=1mm]right:$BN$: $\frac{6}{25}$}] {N} edge from parent node[below,pos=.55,font=\scriptsize,text=teal!70!black,inner sep=1.5pt] {$\frac{2}{5}$}} edge from parent node[above,pos=.55,font=\scriptsize,text=teal!70!black,inner sep=1.5pt] {$\frac{3}{5}$}} child {node[circle,fill=black!80,text=white,inner sep=1.2pt,minimum size=4.2mm,font=\scriptsize\bfseries] {N} child {node[circle,fill=blue!70!black,text=white,inner sep=1.2pt,minimum size=4.2mm,font=\scriptsize\bfseries,label={[font=\scriptsize,xshift=1mm]right:$NB$: $\frac{6}{25}$}] {B} edge from parent node[above,pos=.55,font=\scriptsize,text=teal!70!black,inner sep=1.5pt] {$\frac{3}{5}$}} child {node[circle,fill=black!80,text=white,inner sep=1.2pt,minimum size=4.2mm,font=\scriptsize\bfseries,label={[font=\scriptsize,xshift=1mm]right:$NN$: $\frac{4}{25}$}] {N} edge from parent node[below,pos=.55,font=\scriptsize,text=teal!70!black,inner sep=1.5pt] {$\frac{2}{5}$}} edge from parent node[below,pos=.55,font=\scriptsize,text=teal!70!black,inner sep=1.5pt] {$\frac{2}{5}$}};
\end{tikzpicture}\end{center}a) $\frac35\cdot\frac35=\res{\frac9{25}}$. \ b) $\frac25\cdot\frac25=\res{\frac4{25}}$. \ c) $\frac6{25}+\frac6{25}=\res{\frac{12}{25}}$.

\ex{1} \par\smallskip\begin{center}\begin{tikzpicture}[grow'=right,level distance=17mm,level 1/.style={sibling distance=16mm},level 2/.style={sibling distance=8mm},every node/.style={font=\small\bfseries},edge from parent/.style={draw=muted,thick}]
\node[circle,fill=navy,inner sep=1.2pt] {} child {node[circle,fill=blue!70!black,text=white,inner sep=1.2pt,minimum size=4.2mm,font=\scriptsize\bfseries] {B} child {node[circle,fill=blue!70!black,text=white,inner sep=1.2pt,minimum size=4.2mm,font=\scriptsize\bfseries,label={[font=\scriptsize,xshift=1mm]right:$BB$: $\frac{3}{10}$}] {B} edge from parent node[above,pos=.55,font=\scriptsize,text=teal!70!black,inner sep=1.5pt] {$\frac{2}{4}$}} child {node[circle,fill=black!80,text=white,inner sep=1.2pt,minimum size=4.2mm,font=\scriptsize\bfseries,label={[font=\scriptsize,xshift=1mm]right:$BN$: $\frac{3}{10}$}] {N} edge from parent node[below,pos=.55,font=\scriptsize,text=teal!70!black,inner sep=1.5pt] {$\frac{2}{4}$}} edge from parent node[above,pos=.55,font=\scriptsize,text=teal!70!black,inner sep=1.5pt] {$\frac{3}{5}$}} child {node[circle,fill=black!80,text=white,inner sep=1.2pt,minimum size=4.2mm,font=\scriptsize\bfseries] {N} child {node[circle,fill=blue!70!black,text=white,inner sep=1.2pt,minimum size=4.2mm,font=\scriptsize\bfseries,label={[font=\scriptsize,xshift=1mm]right:$NB$: $\frac{3}{10}$}] {B} edge from parent node[above,pos=.55,font=\scriptsize,text=teal!70!black,inner sep=1.5pt] {$\frac{3}{4}$}} child {node[circle,fill=black!80,text=white,inner sep=1.2pt,minimum size=4.2mm,font=\scriptsize\bfseries,label={[font=\scriptsize,xshift=1mm]right:$NN$: $\frac{1}{10}$}] {N} edge from parent node[below,pos=.55,font=\scriptsize,text=teal!70!black,inner sep=1.5pt] {$\frac{1}{4}$}} edge from parent node[below,pos=.55,font=\scriptsize,text=teal!70!black,inner sep=1.5pt] {$\frac{2}{5}$}};
\end{tikzpicture}\end{center}a) $\frac35\cdot\frac24=\res{\frac3{10}}$. \ b) $\frac25\cdot\frac14=\res{\frac1{10}}$. \ c) $\frac3{10}+\frac3{10}=\res{\frac35}$. \ (Suman 1 \checkmark)

\ex{2} a) $\frac4{40}\cdot\frac3{39}=\frac{12}{1560}=\res{\frac1{130}}$. \ b) $\frac{36}{40}\cdot\frac{35}{39}=\frac{1260}{1560}=\res{\frac{21}{26}}\approx0{,}81$.

\ex{2} a) $\frac5{10}\cdot\frac49=\frac{20}{90}=\res{\frac29}$. \ b) $\frac{20}{90}+\frac3{10}\cdot\frac29+\frac2{10}\cdot\frac19=\frac{20+6+2}{90}=\frac{28}{90}=\res{\frac{14}{45}}$.

\ex{2} \par\smallskip\begin{center}{\small\renewcommand{\arraystretch}{1.2}\begin{tabular}{l|ccc}
\rowcolor{soft}  & Deporte federado & No federado & Total\\\hline
Chicos & 70 & 110 & 180\\
Chicas & 50 & 170 & 220\\
Total & 120 & 280 & 400\\
\end{tabular}}\end{center}a) $\frac{120}{400}=\res{\frac3{10}}$. \ b) $\frac{50}{400}=\res{\frac18}$. \ c) $\frac{220+120-50}{400}=\frac{290}{400}=\res{\frac{29}{40}}$. \ d) Solo contamos las 220 chicas: $\frac{50}{220}=\res{\frac5{22}}$.

\ex{2} \par\smallskip\begin{center}\begin{tikzpicture}[grow'=right,level distance=18mm,level 1/.style={sibling distance=16mm},level 2/.style={sibling distance=8mm},every node/.style={font=\small\bfseries},edge from parent/.style={draw=muted,thick}]
\node[circle,fill=navy,inner sep=1.2pt] {} child {node[] {I} child {node[circle,fill=red!75!black,text=white,inner sep=1.2pt,minimum size=4.2mm,font=\scriptsize\bfseries,label={[font=\scriptsize,xshift=1mm]right:$\frac13\cdot\frac34=\frac14$}] {R} edge from parent node[above,pos=.55,font=\scriptsize,text=teal!70!black,inner sep=1.5pt] {$\frac34$}} child {node[circle,fill=blue!70!black,text=white,inner sep=1.2pt,minimum size=4.2mm,font=\scriptsize\bfseries] {A} edge from parent node[below,pos=.55,font=\scriptsize,text=teal!70!black,inner sep=1.5pt] {$\frac14$}} edge from parent node[above,pos=.55,font=\scriptsize,text=teal!70!black,inner sep=1.5pt] {$\frac13$}} child {node[] {II} child {node[circle,fill=red!75!black,text=white,inner sep=1.2pt,minimum size=4.2mm,font=\scriptsize\bfseries,label={[font=\scriptsize,xshift=1mm]right:$\frac23\cdot\frac14=\frac16$}] {R} edge from parent node[above,pos=.55,font=\scriptsize,text=teal!70!black,inner sep=1.5pt] {$\frac14$}} child {node[circle,fill=blue!70!black,text=white,inner sep=1.2pt,minimum size=4.2mm,font=\scriptsize\bfseries] {A} edge from parent node[below,pos=.55,font=\scriptsize,text=teal!70!black,inner sep=1.5pt] {$\frac34$}} edge from parent node[below,pos=.55,font=\scriptsize,text=teal!70!black,inner sep=1.5pt] {$\frac23$}};
\end{tikzpicture}\end{center}$P(R)=\frac13\cdot\frac34+\frac23\cdot\frac14=\frac14+\frac16=\res{\frac5{12}}$.

\ex{3} Sumamos los caminos con alguna defectuosa ($D$); $B$ es una bombilla buena: $DD$: $\frac3{10}\cdot\frac29=\frac6{90}$; \ $DB$: $\frac3{10}\cdot\frac79=\frac{21}{90}$; \ $BD$: $\frac7{10}\cdot\frac39=\frac{21}{90}$. \ $P=\frac{48}{90}=\res{\frac8{15}}$.

\ex{3} $\frac46\cdot\frac35\cdot\frac24=\frac{24}{120}=\res{\frac15}$.

\ex{3} a) $\res{\frac12}$. \ b) Ana tiene que sacar cruz y Biel cara: $\frac12\cdot\frac12=\res{\frac14}$. \ c) $\frac12+\frac14=\res{\frac34}$.


\end{document}
